% BEGIN GENERATED ARTIFACT: def:truth-table-row-conjunction-propositional-logic
\begin{tcolorbox}[colback=propbox, colframe=propborder, arc=2pt, left=6pt, right=6pt, top=4pt, bottom=4pt, title={\small\textbf{Definition (Truth-Table Row Conjunction)}}]
\begin{definition}[Truth-Table Row Conjunction]
\label{def:truth-table-row-conjunction-propositional-logic}
Fix a finite variable list \(P_1,\ldots,P_n\) and a row \(r\) of the
\hyperref[def:truth-table-propositional-logic]{truth table} for those
variables. The \emph{truth-table row conjunction} for \(r\) is the conjunction
\[
L_1\land L_2\land\cdots\land L_n,
\]
where
\[
L_i=
\begin{cases}
P_i, & \text{if row }r\text{ makes }P_i\text{ true},\\
\neg P_i, & \text{if row }r\text{ makes }P_i\text{ false}.
\end{cases}
\]
\end{definition}
\end{tcolorbox}

\begin{remark*}[Interpretation]
The row conjunction names exactly one truth-table row. It is true on that row
and false on every row that disagrees with it on at least one listed variable.
\end{remark*}

\begin{dependencies}
\begin{itemize}
  \item \hyperref[def:truth-table-propositional-logic]{Truth Table}
  \item \hyperref[def:literal-propositional-logic]{Literal}
  \item \hyperref[def:formula-variable-set-propositional-logic]{Variable Set of a Formula}
\end{itemize}
\end{dependencies}
% END GENERATED ARTIFACT: def:truth-table-row-conjunction-propositional-logic

% BEGIN GENERATED ARTIFACT: def:truth-table-row-clause-propositional-logic
\begin{tcolorbox}[colback=propbox, colframe=propborder, arc=2pt, left=6pt, right=6pt, top=4pt, bottom=4pt, title={\small\textbf{Definition (Truth-Table Row Clause)}}]
\begin{definition}[Truth-Table Row Clause]
\label{def:truth-table-row-clause-propositional-logic}
Fix a finite variable list \(P_1,\ldots,P_n\) and a row \(r\) of the
\hyperref[def:truth-table-propositional-logic]{truth table} for those
variables. The \emph{truth-table row clause} for \(r\) is the clause
\[
M_1\lor M_2\lor\cdots\lor M_n,
\]
where
\[
M_i=
\begin{cases}
P_i, & \text{if row }r\text{ makes }P_i\text{ false},\\
\neg P_i, & \text{if row }r\text{ makes }P_i\text{ true}.
\end{cases}
\]
\end{definition}
\end{tcolorbox}

\begin{remark*}[Interpretation]
The row clause excludes exactly one truth-table row. It is false on row \(r\)
and true on every row that disagrees with \(r\) on at least one listed
variable.
\end{remark*}

\begin{dependencies}
\begin{itemize}
  \item \hyperref[def:truth-table-propositional-logic]{Truth Table}
  \item \hyperref[def:clause-propositional-logic]{Clause}
  \item \hyperref[def:formula-variable-set-propositional-logic]{Variable Set of a Formula}
\end{itemize}
\end{dependencies}
% END GENERATED ARTIFACT: def:truth-table-row-clause-propositional-logic

% BEGIN GENERATED ARTIFACT: def:canonical-disjunctive-normal-form-propositional-logic
\begin{tcolorbox}[colback=propbox, colframe=propborder, arc=2pt, left=6pt, right=6pt, top=4pt, bottom=4pt, title={\small\textbf{Definition (Canonical Disjunctive Normal Form)}}]
\begin{definition}[Canonical Disjunctive Normal Form]
\label{def:canonical-disjunctive-normal-form-propositional-logic}
Let \(\varphi\in\WFF\), and fix an ordering \(P_1,\ldots,P_n\) of the finite
variable set of \(\varphi\). The \emph{canonical disjunctive normal form} of
\(\varphi\) relative to this ordering is the disjunction of the
\hyperref[def:truth-table-row-conjunction-propositional-logic]{truth-table row conjunctions}
corresponding to the rows on which \(\varphi\) is true.
\end{definition}
\end{tcolorbox}

\begin{remark*}[Special case]
If \(\varphi\) is a contradiction, then it has no true rows. Its canonical DNF
is the empty disjunction, represented by the
\hyperref[def:empty-clause-propositional-logic]{empty clause}.
\end{remark*}

\begin{dependencies}
\begin{itemize}
  \item \hyperref[def:disjunctive-normal-form-propositional-logic]{Disjunctive Normal Form}
  \item \hyperref[def:truth-table-row-conjunction-propositional-logic]{Truth-Table Row Conjunction}
  \item \hyperref[def:truth-table-propositional-logic]{Truth Table}
  \item \hyperref[def:empty-clause-propositional-logic]{Empty Clause}
\end{itemize}
\end{dependencies}
% END GENERATED ARTIFACT: def:canonical-disjunctive-normal-form-propositional-logic

% BEGIN GENERATED ARTIFACT: def:canonical-conjunctive-normal-form-propositional-logic
\begin{tcolorbox}[colback=propbox, colframe=propborder, arc=2pt, left=6pt, right=6pt, top=4pt, bottom=4pt, title={\small\textbf{Definition (Canonical Conjunctive Normal Form)}}]
\begin{definition}[Canonical Conjunctive Normal Form]
\label{def:canonical-conjunctive-normal-form-propositional-logic}
Let \(\varphi\in\WFF\), and fix an ordering \(P_1,\ldots,P_n\) of the finite
variable set of \(\varphi\). The \emph{canonical conjunctive normal form} of
\(\varphi\) relative to this ordering is the conjunction of the
\hyperref[def:truth-table-row-clause-propositional-logic]{truth-table row clauses}
corresponding to the rows on which \(\varphi\) is false.
\end{definition}
\end{tcolorbox}

\begin{remark*}[Special case]
If \(\varphi\) is a tautology, then it has no false rows. Its canonical CNF is
the empty conjunction, represented by the
\hyperref[def:empty-conjunction-propositional-logic]{empty conjunction}.
\end{remark*}

\begin{dependencies}
\begin{itemize}
  \item \hyperref[def:conjunctive-normal-form-propositional-logic]{Conjunctive Normal Form}
  \item \hyperref[def:truth-table-row-clause-propositional-logic]{Truth-Table Row Clause}
  \item \hyperref[def:truth-table-propositional-logic]{Truth Table}
  \item \hyperref[def:empty-conjunction-propositional-logic]{Empty Conjunction}
\end{itemize}
\end{dependencies}
% END GENERATED ARTIFACT: def:canonical-conjunctive-normal-form-propositional-logic

% BEGIN GENERATED ARTIFACT: thm:canonical-disjunctive-normal-form-propositional-logic
\begin{tcolorbox}[colback=thmbox, colframe=thmborder, arc=2pt, left=6pt, right=6pt, top=4pt, bottom=4pt, title={\small\textbf{Theorem (Canonical Disjunctive Normal Form)}}]
\begin{theorem}[Canonical Disjunctive Normal Form Theorem]
\label{thm:canonical-disjunctive-normal-form-propositional-logic}
Every formula \(\varphi\in\WFF\) is
\hyperref[def:logical-equivalence-propositional-logic]{logically equivalent}
to its
\hyperref[def:canonical-disjunctive-normal-form-propositional-logic]{canonical disjunctive normal form}
relative to a fixed ordering of its finitely many variables.
\[
\varphi\equiv \operatorname{CDNF}_{P_1,\ldots,P_n}(\varphi).
\]
\noindent\hyperref[prf:canonical-disjunctive-normal-form-propositional-logic]{Go to proof.}
\end{theorem}
\end{tcolorbox}

\begin{remark*}[Interpretation]
Canonical DNF rebuilds the formula from exactly the truth-table rows where the
formula is true. Each row conjunction recognizes one true row, and the final
disjunction accepts any of those rows.
\end{remark*}

\begin{dependencies}
\begin{itemize}
  \item \hyperref[def:canonical-disjunctive-normal-form-propositional-logic]{Canonical Disjunctive Normal Form}
  \item \hyperref[thm:finite-truth-table-propositional-logic]{Finite Truth-Table Theorem}
  \item \hyperref[lem:coincidence-lemma-propositional-logic]{Coincidence Lemma for Propositional Logic}
  \item \hyperref[def:logical-equivalence-propositional-logic]{Logical Equivalence}
\end{itemize}
\end{dependencies}
% END GENERATED ARTIFACT: thm:canonical-disjunctive-normal-form-propositional-logic

% BEGIN GENERATED ARTIFACT: thm:canonical-conjunctive-normal-form-propositional-logic
\begin{tcolorbox}[colback=thmbox, colframe=thmborder, arc=2pt, left=6pt, right=6pt, top=4pt, bottom=4pt, title={\small\textbf{Theorem (Canonical Conjunctive Normal Form)}}]
\begin{theorem}[Canonical Conjunctive Normal Form Theorem]
\label{thm:canonical-conjunctive-normal-form-propositional-logic}
Every formula \(\varphi\in\WFF\) is
\hyperref[def:logical-equivalence-propositional-logic]{logically equivalent}
to its
\hyperref[def:canonical-conjunctive-normal-form-propositional-logic]{canonical conjunctive normal form}
relative to a fixed ordering of its finitely many variables.
\[
\varphi\equiv \operatorname{CCNF}_{P_1,\ldots,P_n}(\varphi).
\]
\noindent\hyperref[prf:canonical-conjunctive-normal-form-propositional-logic]{Go to proof.}
\end{theorem}
\end{tcolorbox}

\begin{remark*}[Interpretation]
Canonical CNF rebuilds the formula by excluding exactly the truth-table rows
where the formula is false. Each row clause rejects one false row, and the
final conjunction requires all false rows to be rejected.
\end{remark*}

\begin{dependencies}
\begin{itemize}
  \item \hyperref[def:canonical-conjunctive-normal-form-propositional-logic]{Canonical Conjunctive Normal Form}
  \item \hyperref[thm:finite-truth-table-propositional-logic]{Finite Truth-Table Theorem}
  \item \hyperref[lem:coincidence-lemma-propositional-logic]{Coincidence Lemma for Propositional Logic}
  \item \hyperref[def:logical-equivalence-propositional-logic]{Logical Equivalence}
\end{itemize}
\end{dependencies}
% END GENERATED ARTIFACT: thm:canonical-conjunctive-normal-form-propositional-logic
